\documentclass{article}

% Language setting
% Replace `english' with e.g. `spanish' to change the document language
\usepackage[english]{babel}

% Set page size and margins
% Replace `letterpaper' with`a4paper' for UK/EU standard size
\usepackage[letterpaper,top=2cm,bottom=2cm,left=3cm,right=3cm,marginparwidth=1.75cm]{geometry}

% Useful packages
\usepackage{amsmath}
\usepackage{graphicx}
\usepackage[colorlinks=true, allcolors=blue]{hyperref}
\usepackage{indentfirst}
\usepackage{enumitem}
\usepackage{physics}
\usepackage{gensymb}
\usepackage{accsupp}

\title{Exercício 6  -  Cálculo Numérico para Geociência}


\begin{document}
\maketitle




Nesta atividade buscaremos retornar aos conceitos apresentados durante o curso, para que possamos revisar os conceitos vistos de interpolação.

Para a determinação de como era o clima há milhares de anos, possuímos a técnica de medir a quantidade de carbono em testemunhos de gelo nas regiões com camadas de gelo que possuem milhares de anos. 

Assim, uma equipe de paleoclimatologistas está estudando a transição da última Era do Gelo para o período interglacial atual através de um testemunho de gelo da Antártica. Eles possuem 4 pontos de controle de concentração de $CO_2$ na atmosfera:

\begin{align*}
\begin{tabular}{|c|c|}
    \hline
    Idade (milhar de anos) & quantidade de $CO_2$ (ppm) \\
    \hline
     22 & 180\\
     18 & 190\\
     14 & 230\\
     12 & 260\\
     \hline
\end{tabular}
\end{align*}

Para estimar a quantidade de carbono em 15 mil anos, dado que não conseguimos tomar, vamos ter que construir um interpolador.

\vspace{0.2 cm}
\hline 
\vspace{0.2 cm}

1) Monte o sistema linear (apenas monte, mas sem resolver) que permite achar os coeficientes na base canônica do polinômio cúbico que interpola os dados acima; 

2) Construa o polinômio cúbico seguindo a forma de Lagrange dada em aula (só explicite os polinômios $L_i$, não precisa simplificar);

3) Construa o polinômio cúbico seguindo a forma de Newton, vista em aula. Depois, simplifique-o para a forma canônica;

4) Substitua os índices encontrados em 3) nos coeficientes a serem determinados em 1). E verifique que é realmente uma solução para nosso problema;

5) Com o polinômio encontrado em 3), determine qual a estimativa de gás carbônico há 15 mil anos.
















\vspace{2cm}
\noindent\centering\large{ESCREVA TODOS OS PASSOS COMO APRESENTADO DURANTE O CURSO}


\end{document}